\section{Introduction}\label{sec:introduction}

Regional governments increasingly use public organizations to connect firms, universities, investors, experts, and other partners around innovation projects. At the same time, firms outside major cities can access dense commercial innovation ecosystems in metropolitan areas. These two forms of intermediation need not be pure competitors at the ecosystem level: they may collaborate, exchange participants, and provide complementary services. Yet for a focal project deciding which intermediary will be its primary route to partners, the alternatives can be substitutes. This coexistence raises a strategic question for decentralized public investment. When one region strengthens its own intermediation capacity, does that induce another region to invest more or less?

The answer is not immediate because innovation intermediation is two-sided. A regional investment that attracts or facilitates partners raises the value of the public route to projects. More projects then make the same route more attractive to additional partners. With a shared private alternative, an investment in one region also shifts projects and partners away from the private route. If the private access price is fixed, that weakening of the shared outside option can raise the marginal return to public investment in the other region. The public investments can therefore be strategic complements.

The private alternative, however, need not passively keep its price unchanged. A commercial innovation intermediary can respond to lost demand by cutting its access price. This response partly restores the private route's attractiveness to projects in both regions. The induced repricing channel works against the positive cross-side network channel. The question of this paper is whether it can be strong enough to change the qualitative strategic relationship between the two decentralized public investors.

We develop a two-region model with two public innovation intermediaries and a shared metropolitan private intermediary. Regional governments simultaneously choose facilitation investments. The private intermediary then chooses a project-side access price. Projects select one primary intermediation route, while support-side partners may participate in multiple ecosystems. Project participation and partner participation reinforce one another. The governments maximize regional welfare; the private intermediary maximizes access-fee profit.

Our main analytic result is a local strategic sign reversal. We compare the full game with a nested benchmark that holds the private price fixed at the full-game equilibrium fee. Consider regular interior stationary branches around the zero-network benchmark, with the G branch anchored at a symmetric regular beta-zero full-game stationary state. In the fixed-price benchmark B3, the public cross derivative is exactly zero at $\beta=0$, but its first derivative with respect to $\beta$ is strictly positive. In the full game G, by contrast, the public cross derivative is already strictly negative at that symmetric regular state when $\beta=0$. Continuity and strict second-order conditions therefore imply that there exists $\varepsilon>0$ such that, for every sufficiently small positive $\beta$ on the continuing regular stationary branches,
\[
\BR_i^{B3\prime}>0>\BR_i^{G\prime}.
\]
This is a sufficient local theorem about the strategic geometry of continuing interior stationary branches. It is not a theorem that those branches are globally optimal for every nearby primitive vector, it does not characterize the entire primitive parameter region, and it does not provide an explicit formula for $\varepsilon$.

The economic reason for the reversal is transparent. With the private price fixed, cross-side participation transmits one region's public expansion through the shared outside option and raises the other region's marginal return to facilitation. Once the commercial intermediary is allowed to optimize its price, public expansion also changes private demand and therefore the private pricing policy. The resulting price response feeds back into project routing in both regions. In derivative form, the full-game public cross effect equals the fixed-price effect plus terms generated by the private price policy and its first and second derivatives. When those repricing terms dominate the fixed-price network force, the public game changes locally from strategic complementarity to strategic substitution. The negative follower-price contribution is not itself claimed as novel; the contribution is the qualitative reversal in the same public--public strategic object when the private pricing policy is activated.

We complement the analytic theorem with a separate repaired global-equilibrium witness. The earlier rational vector used in the first manuscript draft admits profitable finite deviations once all routing regimes are solved and therefore is not used as global-equilibrium evidence. The repaired baseline keeps the model architecture unchanged and uses $\beta=0.01$, $\gamma=0.825$, and $\tau=0.35$, with the other primitives unchanged. An implementation-independent all-regime audit reoptimizes the private continuation after public deviations, searches the full public strategy interval, and checks the previously dangerous low-investment and boundary regions from multiple participation starts. It finds symmetric public best responses near $x_G=0.83710$ in G and $x_{B3}=0.82589$ in B3 and reproduces opposite local slopes, approximately $-0.0199$ and $+0.0040$, respectively. This is computational certification of one repaired baseline witness, not a primitive-space theorem or a uniqueness theorem over parameter values.

The matched-price construction is central to interpretation. B3 uses the same on-path private fee chosen in G but suppresses the private price policy response to public deviations. The comparison therefore isolates the repricing-response margin rather than a difference in the level of the private fee. B3 is an identification benchmark, not a literal price-regulation counterfactual.

Welfare is deliberately separated from strategic interaction. The private access fee is a regional project outflow but an aggregate transfer to the private intermediary, so changing the fee does not mechanically change national welfare. At the repaired full-game witness, a marginal increase in one region's public investment creates a positive effect on the other region's reduced welfare that is only partly offset by lower private profit. The resulting local national coordination wedge is positive. This is a baseline-specific marginal under-provision result; it is not a solved first-best comparison and does not imply that strategic substitution is socially desirable.

The institutional interpretation is also intentionally narrow. Public innovation hubs and commercial metropolitan innovation communities in Japan and abroad exhibit the relevant ingredients: facilitation and matching, cross-side communities, paid access for some private organizations, and access by projects outside the metropolitan core. These examples support the plausibility of the primitives but do not establish the model's causal repricing prediction. In particular, project single-homing is a focal-project primary-route abstraction, private profit maximization is a reduced-form representation of a selected commercial operator class, and the repaired value of the remote-public access friction is constructive rather than empirically calibrated.

Our contribution is not the discovery of network complementarity, endogenous follower pricing, mixed public/private competition, public-investment strategic substitution, or two-sided platforms. Each of those components has close precedents. Sequential-investment and fiscal-competition literatures already show that downstream strategic responses can alter upstream investment incentives, while platform work studies pricing, network feedback, and public/private competition. The narrower contribution is the matched-price comparison for the same two decentralized public investment choices: the local public--public slope is positive when the shared private price response is switched off at the matched full-game fee and negative when that private price setter is activated. The local theorem and the repaired all-regime computational witness support that proposition without converting the familiar component mechanisms into novelty claims.

The rest of the paper proceeds as follows. Section \ref{sec:model} presents the model. Section \ref{sec:equilibrium} characterizes participation, pricing, and the nested benchmarks. Section \ref{sec:mainresult} develops the mechanism decomposition, states the analytic local sign-reversal theorem, and reports the repaired computational witness. Section \ref{sec:welfare} develops welfare and coordination. Section \ref{sec:robustness} separates proved local scope from unproved extension directions. Section \ref{sec:institution} discusses institutional interpretation and empirical implications. Section \ref{sec:literature} positions the result in the related literature. Section \ref{sec:discussion} discusses limitations and scope, and Section \ref{sec:conclusion} concludes.
